\documentclass[11pt]{article}
\usepackage[margin=1in]{geometry}
\usepackage{booktabs}
\usepackage{hyperref}
\usepackage{microtype}

\title{Probability Calibration Under Class Imbalance: A Reproducible Study on UCI Bank Marketing}
\author{Devis Wawan Saputra}
\date{}

\begin{document}
\maketitle

\begin{abstract}
Probability calibration can appear convincing on a single split while remaining sensitive to the split, calibration procedure, metric definition, and feature availability. This reproducible study compares a class-prior baseline with uncalibrated, sigmoid-calibrated, and isotonic-calibrated logistic regression on the UCI Bank Marketing dataset. The protocol uses a frozen external data artifact, repeated stratified holdouts, calibration-fold and ECE-bin sensitivity analyses, an operational feature ablation, error analysis, and descriptive paired uncertainty summaries.
\end{abstract}

\section{Research question}
How do sigmoid and isotonic post-hoc calibration alter logistic-regression probability quality and discrimination on UCI Bank Marketing, and how stable are those changes across repeated splits and sensitivity checks?

\section{Background}
Well-ranked classifiers are not necessarily well calibrated. Classical probability scoring and later empirical calibration work motivate evaluating probability quality separately from discrimination \cite{brier1950,niculescu2005}. Modern calibration literature also emphasizes that confidence estimates can require explicit post-hoc correction and careful evaluation \cite{guo2017}.

\section{Data}
The study uses UCI Bank Marketing, dataset 222 \cite{uciBankMarketing}, derived from direct-marketing campaigns of a Portuguese banking institution and associated with the study by Moro, Rita, and Cortez \cite{moro2014telemarketing}. The executable pipeline extracts \texttt{bank-full.csv} from the official UCI archive. The frozen protocol accepts only the validated source bytes whose SHA-256 is documented in \texttt{DATA.md}; raw data are not committed to the repository.

\section{Methods}
\subsection{Preprocessing}
Numeric predictors are median-imputed and standardized. Categorical predictors are mode-imputed and one-hot encoded with unknown-category tolerance. All preprocessing is fitted inside the training pipeline.

\subsection{Model conditions}
The comparison contains a class-prior baseline, uncalibrated logistic regression, five-fold sigmoid calibration, and five-fold isotonic calibration. Calibration is learned only within the training partition through cross-validation.

\subsection{Evaluation protocol}
The primary split is a stratified 80/20 holdout with seed 42. Robustness uses seeds 13, 29, 42, 73, and 101. The study reports ROC-AUC, average precision, Brier score, log loss, ECE-10, and threshold accuracy. Calibration-fold sensitivity uses 3, 5, and 10 folds; ECE sensitivity uses 5, 10, and 20 equal-width bins.

\subsection{Operational ablation}
Call duration is known only after a call is underway or completed. The primary experiment is therefore rerun without duration to separate a methodological calibration comparison from a pre-contact operational interpretation.

\subsection{Uncertainty and error analysis}
Repeated holdouts reuse observations and are not treated as independent replications. The repository reports descriptive paired split-level bootstrap intervals without inferential p-values. The primary split also records confusion counts, error rate, high-confidence errors, and mean probabilities within false-positive and false-negative subsets.

\section{Generated empirical results}
This section is generated by \texttt{src/run\_experiment.py}; numerical values should not be hand-edited.

Dataset: UCI Bank Marketing (dataset 222), $n=45,211$, 16 predictors, positive prevalence 0.1170.

\begin{table}[htbp]
\centering
\small
\begin{tabular}{lrrrrr}
\toprule
Model & ROC-AUC & AP & Brier $\downarrow$ & Log loss $\downarrow$ & ECE-10 $\downarrow$ \\
\midrule
Class-prior baseline & 0.5000 & 0.1170 & 0.1033 & 0.3609 & 0.0000 \\
Logistic, uncalibrated & 0.9057 & 0.5375 & 0.0720 & 0.2445 & 0.0311 \\
Logistic + sigmoid & 0.9058 & 0.5376 & 0.0720 & 0.2443 & 0.0318 \\
Logistic + isotonic & 0.9058 & 0.5385 & 0.0693 & 0.2269 & 0.0074 \\
\bottomrule
\end{tabular}
\caption{Mean performance across five fixed stratified 80/20 holdouts. Full standard deviations, sensitivity analyses, and error diagnostics are stored in the machine-readable results.}
\label{tab:main-results}
\end{table}

\subsection{Paired robustness summaries}
Across the five fixed splits, isotonic calibration changed Brier score by -0.0027 on average and log loss by -0.0176 relative to uncalibrated logistic regression. These are descriptive robustness summaries rather than independent-replication inference.
Sigmoid calibration produced a much smaller mean Brier change of -0.0000. Exact paired bootstrap intervals are recorded in \texttt{results/metrics.json} and \texttt{paper/results.md}.

\subsection{Operational feature ablation}
The frozen protocol also reruns the primary split after removing call duration, a variable unavailable before a marketing call is completed. The full ablation table is stored in \texttt{results/metrics.json}.

\paragraph{Data identity.}
The exact validated \texttt{bank-full.csv} SHA-256 is \texttt{d1513ec63b385506f7cfce9f2c5caa9fe99e7ba4e8c3fa264b3aaf0f849ed32d}.


\section{Validity and responsible-use boundary}
This is a methodological benchmark on historical marketing data. It does not establish causal effects, present-day population validity, fairness, cross-institution transportability, or suitability for consequential banking decisions. ECE depends on binning, repeated holdouts are dependent, and the dataset may contain repeated contacts with clients. The duration ablation is therefore central to interpreting the study responsibly.

\section{Reproducibility}
The executable protocol is \texttt{src/run\_experiment.py}. Generated evidence is stored in \texttt{results/metrics.json}, \texttt{results/repeated\_runs.csv}, \texttt{results/summary.md}, \texttt{results/figures/}, \texttt{paper/results.md}, and \texttt{paper/results.tex}. Tests and GitHub Actions separate offline code QA from the networked empirical regeneration workflow.

\bibliographystyle{plain}
\bibliography{references}
\end{document}
